\begin{soln}
    We can write an algorithm that takes time $\Theta(n)$. To compute
    $L(\theta)$, we loop through each of the data points, keeping track of the
    loss incurred. If the data point $x$ is red, but $\leq \theta$, we
    increment the loss by one. If $x$ is blue, but bigger than $\theta$, we
    increment the loss by one. Each check takes constant time and there are
    $\Theta(n)$ points, so the function takes $\Theta(n)$ time in total.

    \inputminted{python}{\solutiondir/compute_ell.py}

\end{soln}
